% Laporan praktikum — templat DSWorkbench
% Kompilasi: DSWorkbench: Menulis — Kompilasi (hasil di build/main.pdf)
\documentclass[12pt,a4paper]{article}

\usepackage[utf8]{inputenc}
\usepackage[T1]{fontenc}
\usepackage[indonesian]{babel}
\usepackage[left=3cm,right=2.5cm,top=2.5cm,bottom=2.5cm]{geometry}
\usepackage{amsmath}
\usepackage{graphicx}
\usepackage{booktabs}
\usepackage{float}
\usepackage{listings}
\usepackage[hidelinks]{hyperref}

\graphicspath{{gambar/}}
\lstset{basicstyle=\ttfamily\small, breaklines=true, frame=single, numbers=left, numberstyle=\tiny}

% Identitas: ganti nilai di bawah ini.
\newcommand{\matakuliah}{Nama Mata Kuliah}
\newcommand{\modul}{Modul 1: Judul Modul}
\newcommand{\namamahasiswa}{Nama Mahasiswa}
\newcommand{\nim}{1xx45xxxx}
\newcommand{\kelas}{RA}
\newcommand{\asisten}{Nama Asisten}

\begin{document}

\begin{center}
	{\Large\bfseries Laporan Praktikum \matakuliah}\\[1ex]
	{\large \modul}\\[3ex]
	\begin{tabular}{ll}
		Nama    & : \namamahasiswa \\
		NIM     & : \nim \\
		Kelas   & : \kelas \\
		Asisten & : \asisten \\
		Tanggal & : \today \\
	\end{tabular}
\end{center}

\section{Tujuan}
Tuliskan tujuan praktikum dalam butir yang dapat diperiksa.
\begin{enumerate}
	\item Tujuan pertama.
	\item Tujuan kedua.
\end{enumerate}

\section{Dasar Teori}
Ringkas konsep yang dipakai, dengan kata-kata Anda sendiri. Rumus ditulis seperti
berikut dan dirujuk dengan nomornya, misalnya Persamaan~\ref{eq:rata}.
\begin{equation}
	\bar{x} = \frac{1}{n} \sum_{i=1}^{n} x_i
	\label{eq:rata}
\end{equation}

\section{Alat dan Data}
Sebutkan perangkat lunak beserta versinya dan data yang dipakai (sumber, jumlah baris,
kolom penting).

\section{Langkah Kerja}
Jelaskan langkah yang Anda lakukan secara berurutan. Potongan kode ditulis seperti ini:
\begin{lstlisting}[language=Python, caption={Membaca data}]
import pandas as pd

data = pd.read_csv("data.csv")
print(data.describe())
\end{lstlisting}

\section{Hasil dan Pembahasan}
Sajikan hasil dalam tabel atau gambar, lalu bahas artinya. Contoh tabel ada pada
Tabel~\ref{tab:ringkasan}.
\begin{table}[H]
	\centering
	\caption{Ringkasan hasil}
	\label{tab:ringkasan}
	\begin{tabular}{lrr}
		\toprule
		Variabel & Rata-rata & Simpangan baku \\
		\midrule
		Tinggi   & 165{,}2   & 7{,}4 \\
		Berat    & 58{,}9    & 9{,}1 \\
		\bottomrule
	\end{tabular}
\end{table}

Untuk gambar, simpan berkasnya di folder \texttt{gambar/}, lalu hapus tanda komentar
pada baris berikut.
% \begin{figure}[H]
% 	\centering
% 	\includegraphics[width=0.7\textwidth]{nama-gambar}
% 	\caption{Keterangan gambar}
% 	\label{fig:contoh}
% \end{figure}

\section{Kesimpulan}
Jawab tiap tujuan pada Bagian~1 dengan satu atau dua kalimat berdasarkan hasil.

\section*{Daftar Pustaka}
\begin{enumerate}
	\item Nama Penulis. (Tahun). \textit{Judul buku atau artikel}. Penerbit.
\end{enumerate}

\end{document}
